% !TeX root = ../paper.tex

Touch is the sense that current consumer VR serves worst. A player who reaches for a virtual object and grasps empty air is reminded, at exactly the wrong moment, that none of it is real. We built \emph{Tangible Horror}, a room-scale horror experience for the Meta Quest~3 in which a player wakes up locked inside a cabin and must recover four hidden coins to escape. Each hiding place doubles as a small ordeal: a bowl filled with eyes, a bucket of maggots, a leaking pipe that sprays blood, and cobwebs behind a closet. In the haptic condition, every one of these ordeals has a physical counterpart the player actually touches---rubber balls, gummy worms, water sprayed from a bottle, and synthetic spider webbing---aligned with the virtual scene through spatial anchors and marker-based tracking. In a within-subject study with seven participants (five of whom played both conditions in randomized order), each playthrough was followed by the Igroup Presence Questionnaire (IPQ). The haptic condition scored markedly higher on every IPQ subscale---general presence, spatial presence, involvement, and experienced realism---and every participant who played both versions rated the haptic one higher on all four. Given the small sample, we report these differences descriptively, but their size and consistency are hard to dismiss. We describe the design of the four proxy challenges, the technical setup, and what our results suggest about low-cost passive haptics as a tool for presence in emotionally charged VR.
